\documentclass[11pt,a4paper]{article}

\usepackage[T1]{fontenc}
\usepackage[utf8]{inputenc}
\usepackage[brazil]{babel}
\usepackage{amsmath,amssymb}
\usepackage[margin=2.5cm]{geometry}
\usepackage{enumitem}
\usepackage{tikz}
\usetikzlibrary{arrows.meta,positioning}

\setlist[enumerate]{leftmargin=*}

\newcommand{\Z}{\mathbb{Z}}
\newcommand{\Enc}{\mathsf{Enc}}
\newcommand{\Dec}{\mathsf{Dec}}
\newcommand{\Gen}{\mathsf{Gen}}
\newcommand{\negl}{\mathsf{negl}}
\newcommand{\Priv}{\mathrm{Priv}}
\newcommand{\PubK}{\mathrm{PubK}}
\newcommand{\PubKeav}[3]{\mathsf{PubK}^{\mathrm{eav}}_{#1,#2}(#3)}
\usepackage{titlesec}

\titleformat*{\section}{\normalsize\bfseries}
\titleformat*{\subsection}{\small\bfseries}

\title{Lista 3 \\ EDUARDO KALUF - GRR20241770}
\author{}
\date{}

\begin{document}
    \maketitle

%======================================================================
    \section*{Questão 1 --- ``Seja $\Z_p^* = \{1, 2, \ldots, p-1\}$, onde $p$ é um
    número primo. Usando os resultados vistos em sala, prove que $\Z_p^*$ forma um
    grupo com a operação de multiplicação módulo $p$.''}
%======================================================================

    \paragraph{Fechamento.}
    Para todos $a, b \in \{1, 2, \ldots, p-1\}$ sabemos que o resultado
    de $(a \cdot b) \bmod p$ pertencerá ao grupo. Isso vale pois $p$ é primo e nenhum elemento
    do grupo é divisível por $p$, afinal $p$ é maior que qualquer elemento do grupo. Sendo assim,
    pelo Lema de Euclides, se nem $a$ nem $b$ são divisíveis por $p$, o produto $a \cdot b$
    também não será, ou seja, $(a \cdot b) \bmod p \neq 0$ para quaisquer $a, b \in G$.
    Por último, o resto da divisão por $p$ sempre pertence a $\{0, 1, \ldots, p-1\}$, logo
    $(a \cdot b) \bmod p \in \{1, \ldots, p-1\}$.
    Dessa forma, para todo $g, h \in G$, temos que $(g \cdot h) \bmod p \in G$ também.

    \paragraph{Identidade.} Para qualquer $g \in G$, o elemento $e = 1 \in G$ é definido como o elemento
    identidade, pois para todo $g \in G$ temos que:
    \[
        (g \cdot 1) \bmod p = g = (1 \cdot g) \bmod p.
    \]

    \paragraph{Inverso.}
    Para qualquer $g \in G$ temos que $\gcd(g, p) = 1$, pois $p > g$ e $p$ é primo.
    Isso garante que o número possua um inverso multiplicativo módulo $p$
    pela identidade de Bézout.
    Bézout diz que se temos que $\gcd(g, p) = 1$, então existem $x, y \in \Z$ tais que:

    \[
        g \cdot x + p \cdot y = 1.
    \]

    Tirando o módulo de $p$, o termo $p \cdot y$ se anula, logo:

    \[
        (g \cdot x) \bmod p = 1,
    \]

    Definimos $g^{-1} = x \bmod p$.
    Como $(g \cdot g^{-1}) \bmod p = 1 \neq 0$, temos $g^{-1} \neq 0$, logo $g^{-1} \in G$ é o
    inverso multiplicativo de $g$.

    \paragraph{Associatividade.} A associatividade das operações segue de maneira trivial.
    Multiplicar dois inteiros é uma operação associativa e temos a seguinte propriedade:

    \[
        ((x \bmod p) \cdot y) \bmod p = (x \cdot y) \bmod p = (x \cdot (y \bmod p)) \bmod p
    \]

    Sendo assim, para quaisquer $a, b, c \in G$ temos que:

    \[
        \bigl(((a \cdot b) \bmod p) \cdot c\bigr) \bmod p
        = (a \cdot b \cdot c) \bmod p
        = \bigl(a \cdot ((b \cdot c) \bmod p)\bigr) \bmod p.
    \]


    Sendo assim, $\Z_p^*$ atende todas as definições e provamos que ele é um grupo.

    \clearpage

%======================================================================
    \section*{Questão 2 --- ``(Exercício 10.3 do livro, adaptado). Descreva um ataque
    `Man-in-the-middle' para o protocolo de troca de chaves de Diffie-Hellman. No fim
    do protocolo, o adversário deve ter compartilhado uma chave $k_A$ com Alice e uma
    chave $k_B$ com Bob, potencialmente diferente de $k_A$, e Alice e Bob não
    detectam o que ocorreu.''}
%======================================================================

    Alice quer compartilhar uma chave secreta com Bob utilizando o protocolo Diffie-Hellman.
    Dessa forma, Alice deve primeiramente sortear um número aleatório $x \leftarrow \Z_q$ e, em seguida,
    enviar $h_A = g^x$, tal que $g \in G$, para Bob.

    No meio do caminho, a mensagem é interceptada por Maine, que tem controle sobre o canal de comunicação.
    Maine recebe $h_A$ de Alice e o guarda, sorteando $m_1 \leftarrow \Z_q$ e repassando $h_{m_1} = g^{m_1}$ para Bob.

    Bob recebe $h_{m_1}$ de Maine, pensando ter recebido de Alice, e faz o mesmo: sorteia $y \leftarrow \Z_q$
    e envia $h_B = g^y$ para Alice.
    Novamente, Maine intercepta o valor, guarda $h_B$, sorteia $m_2 \leftarrow \Z_q$ e repassa $h_{m_2} = g^{m_2}$ para Alice.

    As chaves são definidas como a exponenciação entre os valores recebidos e seus expoentes secretos:

    \[
        K_A = (h_{m_2})^{x} = g^{x m_2}, \qquad K_B = (h_{m_1})^{y} = g^{y m_1}.
    \]

    Como Maine conhece $m_1$ e $m_2$, ele consegue produzir as mesmas chaves de Alice e de Bob:

    \[
        K_A = (h_A)^{m_2} = g^{x m_2}, \qquad K_B = (h_B)^{m_1} = g^{y m_1}.
    \]

    Assim, Maine possui a chave $K_A$, a mesma de Alice, e a chave $K_B$, a mesma de Bob.

    Essas duas chaves têm probabilidade altíssima de serem diferentes.
    Elas são iguais somente quando os valores sorteados por Maine acabam fazendo com que ambas as contas deem o mesmo resultado.

    Alice e Bob não conseguem detectar o ataque, pois todos os elementos são sorteados de maneira uniforme
    e o protocolo não realiza autenticação das mensagens enviadas.
    Após estabelecer a comunicação, Maine continua interceptando as mensagens, decifrando cada mensagem de
    Alice com $K_A$ e recifrando-a com $K_B$ antes de repassá-la a Bob (e vice-versa), mantendo a
    comunicação normal para ambos os lados honestos.

    \begin{center}
        \begin{tikzpicture}[>=Stealth]
        \node[align=center] at (0,-0.4) {Alice\\[2pt] $x \leftarrow \Z_q$\\[2pt] $h_A = g^x$};
        \node[draw, minimum width=2.6cm, minimum height=2.8cm, label={[below=2pt]north:Maine}] (mal) at (5,-0.6) {};
        \node[font=\small, align=center] at (5,-0.9) {$m_1 \leftarrow \Z_q$\\[2pt] $m_2 \leftarrow \Z_q$};
        \node[align=center] at (10.4,-0.4) {Bob\\[2pt] $y \leftarrow \Z_q$\\[2pt] $h_B = g^y$};

        \draw[->] (1.2,-0.2) -- node[above]{$h_A$}     (3.7,-0.2);
        \draw[<-] (1.2,-1.0) -- node[above]{$h_{m_2}$} (3.7,-1.0);
        \draw[->] (6.3,-0.2) -- node[above]{$h_{m_1}$} (9.4,-0.2);
        \draw[<-] (6.3,-1.0) -- node[above]{$h_B$}     (9.4,-1.0);

        \node at (0,-2.0)    {$K_A = (h_{m_2})^{x}$};
        \node at (10.4,-2.0) {$K_B = (h_{m_1})^{y}$};
        \node[below=4pt of mal, align=center] {$K_A = (h_A)^{m_2}$\\[2pt] $K_B = (h_B)^{m_1}$};
        \end{tikzpicture}
    \end{center}

    \clearpage

%======================================================================
    \section*{Questão 3}
%======================================================================

    \begin{quote}\small\bfseries
        ``Seja $\Pi = (\Gen, \Enc, \Dec)$ uma cifra de chave pública que é CPA-segura com
        espaço de mensagens $\{0,1\}^n$ para parâmetro de segurança $n$. Neste exercício,
        vamos mostrar que `o jeito óbvio' de estender $\Pi$ para mensagens de tamanho
        $\{0,1\}^{2n}$ é também CPA-seguro. Em particular, vamos mostrar que a cifra
        $\Pi' = (\Gen, \Enc', \Dec')$ onde
        $\Enc'_{pk}((m_1, m_2)) = (\Enc_{pk}(m_1), \Enc_{pk}(m_2))$ e
        $\Dec'_{sk}((c_1, c_2)) = (\Dec_{sk}(c_1), \Dec_{sk}(c_2))$ é também CPA-segura.

        Para a demonstração, usaremos o método híbrido, com três experimentos híbridos
        baseados no experimento $\PubKeav{\Pi'}{\mathcal{A}}{n}$ (lembre que uma cifra de
        chave pública EAV-segura é também CPA-segura) definidos a seguir:

        \begin{itemize}
            \item $E_{00}$ é idêntico ao experimento $\PubKeav{\Pi'}{\mathcal{A}}{n}$, porém
            o desafiante, ao receber as mensagens de desafio $(m_{0,0}, m_{0,1})$ e
            $(m_{1,0}, m_{1,1})$, sempre retorna $\Enc'_{pk}(m_{0,0}, m_{0,1})$.
            \item $E_{01}$ é também idêntico ao experimento $\PubKeav{\Pi'}{\mathcal{A}}{n}$,
            porém o desafiante sempre retorna $\Enc'_{pk}(m_{0,0}, m_{1,1})$. Isto é, o
            desafiante produz um texto cifrado que corresponde à primeira metade da
            primeira mensagem e à segunda metade da segunda mensagem.
            \item $E_{11}$ é também idêntico ao experimento $\PubKeav{\Pi'}{\mathcal{A}}{n}$,
            porém o desafiante sempre retorna $\Enc'_{pk}(m_{1,0}, m_{1,1})$.
        \end{itemize}

        Perceba que o experimento $E_{01}$ não corresponde à uma situação que de fato
        possa ocorrer no experimento $\PubKeav{\Pi'}{\mathcal{A}}{n}$. Ainda assim, nada nos
        impede de analisar o comportamento de $\mathcal{A}$ nessa situação (e essa análise
        é útil para nossa demonstração).''
    \end{quote}

    \subsection*{a)}

    \begin{quote}\small\bfseries
        ``Defina $p_{00}$, $p_{01}$ e $p_{11}$ como a probabilidade de que o adversário
        $\mathcal{A}$ responda $b' = 1$ no experimento $E_{00}$, $E_{01}$ e $E_{11}$,
        respectivamente. Argumente que
        \[
            \Pr\!\left[\PubKeav{\Pi'}{\mathcal{A}}{n} = 1\right]
            = \frac{1}{2} + \frac{1}{2} \cdot (p_{11} - p_{00}).\text{''}
        \]
    \end{quote}

    O adversário ganha $E_{00}$ quando responde $0$.\\
    O adversário nunca ganha $E_{01}$.\\
    O adversário ganha $E_{11}$ quando responde $1$.

    Ou seja, o experimento todo pode ser medido como $\frac12$, chance de ganhar
    chutando, mais a situação em que ele vence os desafios, % leitura incerta no manuscrito
    que podem ser descritos em termos de $p_{00}, p_{01}, p_{11}$ assim:
    \[
        \Pr = \frac12 + \frac12\,(p_{11} - p_{00}).
    \]

    \subsection*{b)}

    \begin{quote}\small\bfseries
        ``Prove, a partir da segurança de $\Pi$, que existe uma função negligenciável
        $\negl$ tal que $|p_{01} - p_{00}| \le \negl(n)$.

        \textit{Dica.} Apresente uma redução que transforma um adversário $\mathcal{A}$
        para o qual $|p_{01} - p_{00}| = \varepsilon$ em um adversário $\mathcal{A}'$ que
        possui vantagem similar para quebrar a segurança de $\Pi$.

        Explique, de maneira sucinta, como a sua demonstração pode ser adaptada para
        provar também que existe uma função negligenciável $\negl$ tal que
        $|p_{11} - p_{01}| \le \negl(n)$.''
    \end{quote}

    Podemos utilizar $A$ como subrotina no desafio $\Priv_{A'}$ a fim de quebrar
    $\Pi$, basta que $A'$ responda $1$ quando $A$ responder $1$.

    \begin{center}
        \begin{tikzpicture}[>=Stealth, box/.style={draw, minimum width=1.4cm, minimum height=1.6cm}]
            \node[box, label=above:Desafiante] (d)  at (0,0)   {};
            \node[box, label=above:$A'$]       (ap) at (4,0)   {};
            \node[box, label=above:$A$]        (a)  at (7.5,0) {};
            \draw[->] ([yshift=4mm]d.east)   -- ([yshift=4mm]ap.west);
            \draw[<-] ([yshift=-4mm]d.east)  -- node[below]{$b'$} ([yshift=-4mm]ap.west);
            \draw[->] ([yshift=4mm]ap.east)  -- ([yshift=4mm]a.west);
            \draw[<-] ([yshift=-4mm]ap.east) -- ([yshift=-4mm]a.west);
        \end{tikzpicture}
    \end{center}

    Dessa forma, quebraríamos $\Priv_{A'}$, porém sabemos que
    \[
        \Priv_{A'} \le \frac12 + \negl(n),
    \]
    logo deve ser para $A$ também:
    \begin{align*}
        \Priv_A = \frac12 + \frac12\,(p_{01} - p_{00}) &\le \frac12 + \negl(n)\\
        p_{01} - p_{00} &\le 2\,\negl(n) = \negl(n)\\
        |p_{01} - p_{00}| &\le \negl(n)
    \end{align*}
    O mesmo raciocínio segue para $|p_{11} - p_{01}|$.

    \subsection*{c)}

    \begin{quote}\small\bfseries
        ``Argumente que $\Pi'$ é CPA-segura.

        \textit{Dica.} Use a seguinte propriedade de diferenças modulares:
        \[
            |a - c| = |a - b + b - c| \le |a - b| + |b - c|.\text{''}
        \]
    \end{quote}

    Pela desigualdade triangular:
    \[
        |p_{11} - p_{00}| = |p_{11} - p_{01} + p_{01} - p_{00}|
        \le |p_{11} - p_{01}| + |p_{01} - p_{00}|.
    \]
    Logo:
    \[
        \PubK = \frac12 + \frac12\,(p_{11} - p_{00})
        \le \frac12 + \frac12\bigl(|p_{11} - p_{01}| + |p_{01} - p_{00}|\bigr).
    \]
    Pelo exercício b), temos que
    \begin{align*}
        \PubK &\le \frac12 + \frac12\bigl(\negl(n) + \negl(n)\bigr)\\
        &\le \frac12 + \negl'(n).
    \end{align*}
    Sendo assim, $\Pi'$ é seguro.

    \clearpage

%======================================================================
    \section*{Questão 4}
%======================================================================

    \begin{quote}\small\bfseries
        ``(Exercício 12.4 do livro, adaptado). Considere o seguinte experimento de
        segurança para assinaturas digitais: um adversário probabilístico polinomial
        $\mathcal{A}$ recebe uma chave pública $\langle N, e \rangle$ e uma mensagem
        $m \leftarrow \Z_N^*$ escolhida de maneira uniforme, e produz uma `assinatura'
        $\sigma$. O adversário vence se $\sigma$ é uma assinatura válida para $m$. Perceba
        que o adversário não pode realizar nenhuma consulta de assinatura.

        Assumindo que o problema RSA é difícil, mostre que, para o esquema de assinatura
        RSA básico que vimos em sala obedece ao seguinte: para todo adversário
        probabilístico polinomial $\mathcal{A}$ existe uma função negligenciável $\negl$
        tal que
        \[
            \Pr\!\left[\mathcal{A}(\langle N, e \rangle, m) \text{ vence no experimento}\right]
            \le \negl(n).
        \]
        Onde $n$ é o parâmetro de segurança usado para gerar a chave
        $\langle N, e \rangle$ e a aleatoriedade é sobre as escolha de
        $\langle N, e \rangle$, $m$ e sobre a aleatoriedade de $\mathcal{A}$.''
    \end{quote}

    % Resposta da Questão 4

\end{document}